\chapter{Understanding the Dynamics of Gender Stereotypes in Mathematics}\label{chap:SocialScience}
\section{Introduction}
Currently, the majority of college students are women, as seen in the data analyzed by \textcite{OWID2024}. However, the number of females pursuing careers in STEM (Science, Technology, Engineering, and Mathematics) fields remains low, raising concerns among educators and policymakers.
For instance, research by \textcite{eccles1989} stated that in 1983, only 13.2\% of all undergraduate engineering degrees were earned by women, and only 27\% in physical sciences. In contrast, women's representation in fields such as psychology and education was much higher, at 68\% and 75\%, respectively. The gap becomes even more obvious at the graduate and professional levels, with less than 10\% of master's degrees and less than 6\% of Ph.D. degrees in engineering awarded to women. This inadequate representation leads to a significant loss of talent and diversity in STEM fields, which are vital for fostering innovation and progress.

The issue goes beyond basic aptitude or access, revealing more subtle structural barriers that persist even in an era of increasing educational opportunities. The paradox of women performing well academically overall—often achieving higher scores in science and math during high school—yet choosing to limit their interest in these subjects suggests that more complex factors are involved than mere bias. These trends indicate that the issue may come from intricate psychological processes and reactions to contextual cues, rather than being solely external (\cite{Monteiro2023}).

Understanding preconceptions, especially negative ones, is crucial for grasping this phenomenon. Stereotypes can significantly affect individuals' self-perception, motivation, and performance, particularly when these ideas target their social group, as was seen in the research by \textcite{HEILMAN2012113}. The term "stereotype threat" describes the psychological anxiety that arises from the awareness that one might be judged based on a negative stereotype (\cite{Spencer2016}).

\section{Women's participation and achievement in Mathematics}
The visible differences in participation and achievement between genders in mathematics have changed over time, making this issue dynamic rather than static.
Historically, large-scale studies have often shown a male advantage in math proficiency, particularly noticeable by the end of high school and in more complex cognitive tasks (\cite{LiuWilson2009}). For instance, the early assessment known as Project TALENT (1960) and the International Study of Achievement in Mathematics (1967) pointed out a modest male advantage by the end of high school and in higher-order cognitive tasks. These disparities were often attributed to differences in course-taking patterns, with males more likely to pursue advanced mathematics in secondary education (\cite{Armstrong}).

By the late 1970s, some assessments like the Women in Mathematics Project (1978) and the National Assessment of Educational Progress (NAEP, 1977-78) — that took place in the United States — present a more nuanced picture of gender differences in mathematics. While the participation gap in general mathematics courses—such as General Mathematics, Algebra 1, and Geometry—has decreased, disparities still favor males at higher levels of mathematics and in specific achievement areas, such as problem-solving (\cite{Armstrong}).

In the late 2000s, \textcite{Zel} found that 13-year-old girls often begin high school with mathematical abilities equal to or even better than those of boys. On occasion, they perform better than their male counterparts in computation and spatial visualization. However, by the end of high school, a noticeable change occurs: males in the 12th grade score higher on problem-solving tests, while girls lose their former advantages in computation and spatial skills. While early achievement gaps could be explained by different participation and interest in math and math-related courses, data show that disparities persist even when participation is equivalent. For example, males continue to outperform females on complex word problems and application-based tasks, revealing that the issue extends beyond mere exposure (\cite{Zel}).

Some hypotheses suggest that males have an advantage due to their superior spatial visualization skills (\cite{RamirezUcles2020}); however, the research by (\cite{FennemaTartre1985}) shows no relevant differences in spatial ability between sexes. Instead, the ongoing divergence in achievement indicates that other factors have a role. These factors include: learning opportunities outside of a formal education environment (\cite{Boeren2011}), problem-solving strategies (\cite{RamirezUcles2020}), and psychological traits, such as motivation, self-confidence, self-efficacy, and perseverance (\cite{Rodriguez2020}). These results highlight the importance of social and psychological factors in affecting engagement and achievement in mathematics, particularly at higher levels.

\section{Theoretical Framework for Understanding Gender Differences in Mathematics}
To thoroughly explain the gender disparities in mathematics, it is essential to have a theoretical framework that considers the psychological and social factors influencing students' attitudes and behaviors. Two well-established models that offer a comprehensive understanding of how these elements affect academic decisions and performance are Eccles' Expectancy-Value Theory and Bandura's Social Cognitive Theory.

\subsection{Eccles' Expectancy-Value Theory}
Eccles' model of academic choice (\cite{eccles1989}) highlights two key factors that significantly influence an individual's academic decisions: their assessment of a task and their expectations for success. The central idea of this model is that individuals' expectations, values, and behaviors related to achievement are primarily shaped by their perceptions of reality, rather than by their actual past performances. This means that a student's confidence in mathematics and the importance they assign to the subject are deeply impacted by their self-assessment of ability and how others perceive their success (for example, attributing their achievements to effort rather than innate talent), even when they consistently receive high grades.

The model specifically outlines several mediating factors that affect both expectations and values. These factors include established norms of gender roles, the influence of key socializers such as parents and teachers, perceptions of task difficulty, patterns of causal attribution, and the ways individuals define their successes and failures. Together, these elements significantly shape a person's perspective, which directly influences their level of engagement, effort, and ultimately, their academic and professional choices.

Research consistently shows that boys and girls have significantly different attitudes toward English and mathematics, as well as different perceptions of themselves as learners, particularly after the age of 12. These disparities typically begin in the seventh grade and become more pronounced throughout high school. Even when girls achieve grades equal to or better than boys', they often develop lower self-assessments regarding their mathematical abilities. 

There is a growing gap between girls' confidence in their math skills and their confidence in their English abilities, with girls generally feeling more confident in English. In contrast, this consistent pattern of confidence levels is not usually observed among boys. This suggests that many girls' self-confidence in mathematics declines because of their subjective evaluations of their skills, overshadowing their actual academic achievements.

This pattern is reflected in the way individuals perceive the value of mathematics subjectively. In terms of math assessment, females demonstrate a consistent decline over time regarding their perception of math's value, which includes its usefulness, significance, and their interest in it. In contrast, males’ assessments tend to remain relatively stable. Consequently, as they progress through high school, girls come to value English more than math. 

Applying Eccles' theory, the reason women are less likely to choose advanced math or physical science courses—preferring instead to pursue advanced studies in English, foreign languages, or social sciences—can be partly attributed to this difference in confidence and perceived worth. For these students, abandoning math is often not due to an inherent dislike of the subject; rather, it is a conscious decision to focus on other interests or objectives they find more rewarding. This perspective indicates that academic choices involve a strategic trade-off, as mathematics is viewed as less useful or less aligned with a woman's evolving self-concept compared to other subjects in which she feels more competent and valued.

\subsection{Bandura's Social Cognitive Theory}
Self-efficacy is the key construct in Albert Bandura's Social Cognitive Theory (\cite{Zel,Bandura2001}). It is defined as individuals' judgments of their ability to achieve specific performance goals. These perceptions significantly influence motivation, decision-making, effort, perseverance, and the ability to recover from setbacks.
Individuals rely on four main sources to form their perceptions of self-efficacy.
\begin{itemize}
    \item Authentic Mastery Experiences: These kinds of experiences come from past performances and are considered the most influential source of self-efficacy. While repeated failures, especially in the early stages and at a young age, can greatly diminish a person's sense of efficacy, success in a task can foster a strong belief in one’s abilities.
    \item Vicarious Experiences: These come from observing others perform tasks. The impact of these observations increases when the observer has little experience in the area, or when the models are perceived to be similar to the observer.
    \item Social Persuasions: These comments reflect others' perceptions of an individual's capability to complete a task. People who already have a strong sense of self-efficacy benefit from positive encouragement, which motivates them to make more effort. However, derogatory remarks can damage the self-esteem of those who are already lacking self-assurance.
    \item Physiological Indexes: People assess their abilities based on their mental and physical states, such as stress and anxiety. These conditions can influence self-efficacy, either enhancing or undermining confidence.
\end{itemize}
The construct of self-efficacy was introduced by \textcite{HACKETT1981326} to explain the underrepresentation of women in higher-status, male-dominated fields (\cite{Zel}). They proposed that cultural expectations are internalized and influence career choices, leading women to limit their options due to a lack of self-efficacy in career-related activities. It was discussed that differences in experiences related to self-efficacy sprout from varying gender-role socialization. For example, boys are often exposed to mechanical, scientific, and technical activities at an earlier age, which can enhance their confidence in these areas. Moreover, men are more likely to encounter relatable role models in these fields.

Research indicates significant gender differences in the primary sources from which individuals in STEM fields develop their self-efficacy:
\begin{itemize}
    \item Women in STEM: Social persuasion and vicarious experiences are key sources for developing and maintaining self-efficacy among women in STEM professions. Often, women recall more episodes of supporting others and witnessing their achievements than their achievements. Their beliefs about themselves are significantly shaped by peers, academic environments, families, and professional settings. In particular, in male-dominated fields, women frequently mention the positive impact of influential individuals who have helped them build both confidence and skills. This research suggests that, perhaps even more than individual accomplishments, the social and relational context surrounding women's experiences is crucial for establishing and sustaining confidence in these male-dominated areas.
    \item Men in STEM: Men working in STEM fields often identify mastery experiences as the primary source of their self-efficacy. They emphasize individual accomplishments and stories of self-motivation, frequently referring to innate talent or skill. While they also recognize the impact of social encouragement and vicarious experiences, these factors typically serve as reinforcement rather than the main roots of their confidence. Men often view themselves as passively influenced by these external factors, interpreting vicarious experiences as enhancements to their self-efficacy and providing role models for entering STEM careers. 
\end{itemize}
These different pathways demonstrate that a uniform approach to fostering self-efficacy may not be effective. Recognizing these differences is crucial for designing targeted interventions.
The diverse sources of self-efficacy also reveal different patterns of resilience and identity development:
\begin{itemize}
    \item Resilience of Women in STEM: Women pursuing careers in STEM often view themselves as resilient when facing both social and intellectual challenges. They draw strength from social influences and the vicarious experiences gained through important relationships. The encouragement they receive from others boosts their self-confidence, enabling them to overcome obstacles related to being women in a male-dominated field. To counteract the negative effects of societal pressures, many women choose to ignore harmful social cues or develop effective coping strategies. They face challenges related to their social identity and the psychological stress of navigating their gender identity in professions that may subtly or overtly challenge traditional feminine ideals. According to Gilligan's theory, this "identity cost" requires a strong sense of self-efficacy, which is often fostered through social support (\cite{Gilligan1993}).
    \item Men's Resilience: Men tend to have a strong sense of self-efficacy, which is largely shaped by their mastery experiences. As a result, they are more likely to view adversity as a challenge rather than a harmful obstacle. They typically focus on clarifying their goals instead of doubting their ability to succeed in STEM fields. Additionally, they generally do not face relevant issues with social identity challenges related to their profession.
\end{itemize}
The development of self-efficacy varies between genders, aligning with Erik Erikson's theories (1959, 1968) (\cite{Zel}). Erikson suggested that men's identity is often linked to independent, work-related achievements, operating in an "outer space." In contrast, women's identity construction, as Erikson described and Carol Gilligan later expanded upon (1982), is more closely tied to an "inner space" defined by the quality of their relationships. Gilligan argued that for women, intimacy and identity are interconnected, as they rely on their interactions with others to define and understand themselves, which subsequently influences their behavior. This suggests that for women in STEM, their significant others play a vital role in reinforcing confidence and their perseverance, helping them see themselves as mathematicians and scientists.

\section{The Impact of Stereotype Threat on Women in Mathematics}
Stereotype threat refers to the psychological fear that arises when individuals know they are being viewed through the lens of a negative stereotype (\cite{Beeghly2015}). This phenomenon can significantly impact people's performance, particularly in areas where such stereotypes target their group (\cite{Schmader2008}). As a result, performance may decline, which can unintentionally reinforce the very stereotype that triggered the threat. This effect is often most pronounced in individuals who have a strong identification with the achievement domain, as their success in that area is closely tied to their sense of self-worth (\cite{CROCKER199989}).
Research consistently indicates that women's performance in mathematics suffers when they experience "solo status"—being the only woman in an all-male group. This impairment is explained through a three-step process by \textcite{Beat}:
\begin{itemize}
    \item Numerical Distinctiveness: A woman stands out in a group as the only female, making her stand out. Her gender receives more attention due to her heightened visibility.
    \item Stereotype Threat: The belief that "women aren't suited for math" is an example of negative stereotypes that may arise from individual differences, highlighting gender as a factor. This creates a "threatening intellectual environment" where individuals anticipate being judged more based on group stereotypes than on their abilities.
    \item Performance deficits: Unfavorable preconceptions can create pressure and anxiety, which may directly impact cognitive function and performance on the job.
\end{itemize}
\subsection{Mediating Mechanisms}
Literature also explores two primary mechanisms through which stereotypes affect behavior:
\begin{itemize}
    \item Cognitive Explanation (Stereotype Activation): This perspective suggests that poor performance may come from the activation of stereotypes. According to this theory, when stereotypes are activated, they can disrupt cognitive functions, making it more difficult to build proper solutions and lowering overall performance standards. \textcite{Beat}, found that single women slightly outperformed control groups when completing stereotyped word fragments, indicating some level of stereotype activation. However, this activation was not consistently linked to deficiencies in mathematical abilities. This raises the possibility that the observed performance disadvantages may not be solely attributed to the presence of stereotypes. Instead, a more significant factor contributing to performance issues appears to be the psychological distress associated with being judged based on stereotypes, rather than the mere mental presence of the stereotype itself.
    \item Motivational Explanation (Stereotype Anxiety): This perspective suggests that the primary cause of performance declines is the experience of threat-like events and emotional reactions, such as anxiety. When faced with the threat of stereotypes, single women often become overly focused on how they will be judged based on these stereotypes. This preoccupation diverts their attention from their actual work, ultimately hindering their focus and performance. In \textcite{Beat}, this idea was partially supported, showing that single women exhibited significantly higher levels of stereotype anxiety compared to women who were not solo. This increased anxiety was associated with a decrease in mathematics performance. These results suggest that emotional and psychological distress, arising from fear of being judged based on stereotypes, contributes significantly to impaired performance.
\end{itemize}
\subsection{Ego-protective Response: Identity Bifurcation}
Individuals who experience strong or persistent stereotype threat may engage in various ego-protective behaviors (\cite{SnyderMiene1994}). One common response is disengagement or disidentification with the affected domain, where the person reduces their performance in that area to protect their self-evaluation. However, a more nuanced and less costly approach is "identity bifurcation." This involves maintaining a connection to other valued qualities of one’s social group that are not associated with stereotypes, while selectively distancing oneself from aspects linked to negative perceptions in the threatened area.
Empirical evidence for identity bifurcation was provided by \textcite{pro}. In Study 1, women in a real-world academic setting who had taken a large number of math classes strongly disavowed "feminine characteristics" associated with negative math stereotypes (such as flirtatiousness and plans to have children) more than women who had taken fewer math courses. However, they did not reject characteristics that are weakly associated with negative stereotypes, such as empathy and nurturance. 

In Studies 2 and 3, when stereotype threat was experimentally induced by presenting a scientific article that reported sex differences in math aptitude, math-identified women responded by rejecting strongly associated feminine characteristics (like wearing make-up and being artistic), but not weakly associated ones such as empathy and having a good fashion sense. In contrast, women who had a weak identification with math did not show this pattern of bifurcation. 

It is important to note that in neither study was the abandonment of feminine traits accompanied by an increase in identification with masculine traits.

Identity bifurcation is a response intended to protect the ego, but it comes with a significant hidden cost. While it allows women to navigate challenging environments, it often requires them to sacrifice parts of their identity, which can lead to an unnecessary struggle to fit in. This adaptation, although protective, is a sign of an inequitable environment.

Stereotype threat is a phenomenon that occurs in specific contexts rather than a fixed trait within individuals. Characteristics that some women may disavow, like being artistic or interested in family, are not negative in themselves but can become threatening in math-related situations, such as classrooms or during challenges to mathematical abilities. This indicates that the real issue lies in the interaction between individuals and their environments.
Additionally, scientific reports that reinforce stereotypes about ability can create a "self-fulfilling prophecy." When these stereotypes are emphasized by peers or media, they can negatively impact performance, raising important questions about how research findings are communicated.

\section{Effects of Stereotypes on Mathematics Performance}
As seen in the previous sections, the literature surrounding gender stereotypes in mathematics reveals that these beliefs strongly impact women's participation, achievement, and self-perception within the field. 
As early as the 1980s, extensive surveys indicated systematics disparities in achievemtn and enrollment patterns, demonstrating that women were consistently underrepresented in advanced mathematics courses and tended to score lower on standardized assessments compared to their male counterparts (\cite{Armstrong}).
Such important discrepancies cannot be attribuited solely to differences in ability; rather, they are deeply connected to the societal perception of mathematics as a male-dominated field (\cite{FennemaTartre1985,Plante2009}).
Recent research has clarified the psychological mechanisms by which stereotypes impact mathematics performance. Studies indicate that stereotype threat can hinder performance by increasing anxiety, reducing working memory capacity, and diverting cognitive resources toward self-monitoring rather than focusing on the task at hand (\cite{Schmader2008,Spencer2016}). Furthermore, experimental studies have shown that even merely witnessing another woman being subjected to dominant male behavior in a specifically mathematical context  can trigger stereotype threat and decrease women's performance. This highlights the subtle nature of these effects (\cite{VanLooRydell2014}).

Stereotypes also influence individuals indirectly by shaping their beliefs about self-efficacy and academic motivation.  In mathematics, young girls who are exposed to gender-stereotypical expectations from parents and teachers often develop lower confidence about their own competence. This leads to lower engagement and achievement in the subject (\cite{Jacobs1991,Gunderson2012}). These processes result in cumulative disadvantages: women generally report lower self-efficacy in mathematics and are less likely to pursue advanced studies or careers in mathematics, even if they demonstrate equal or superior performance (\cite{Zel,Beat}).

Cross-cultural evidence highlights these common dynamics. Studies comparing China and the United States show that gender stereotypes impact math performance in both countries, although the size of the gap varies with cultural values and school systems (\cite{Tsui2007,LiuWilson2009}). Additionally, classmates' beliefs can affect individual results: girls surrounded by peers who hold traditional views of math tend to do worse, while boys' performance stays the same or even gets better (\cite{Smith22}).

The long-term effects are clear in participation patterns. While women have made significant strides in tertiary education worldwide, they remain underrepresented in mathematics-intensive fields (\cite{OWID2024,Boeren2011}). Scholars argue that this is not just about individual choices; it's also shaped by stereotypes that frame math as inconsistent with female identity and goals (\cite{eccles1989,Gilligan1993,Monteiro2023}). This can lead to “identity bifurcation,” where women distance themselves from aspects of their identity linked to math to protect their self-esteem (\cite{pro,CROCKER199989}).

So, gender stereotypes in mathematics do not just reflect performance differences, but they actively contribute to creating them. They impact performance through various channels, including direct cognitive interference, long-term effects on self-efficacy and motivation, and social pressures that deter participation. Although the effects vary by culture, development, and context, their cumulative impact is significant, perpetuating gender gaps in mathematics achievement and representation. Addressing these stereotypes is crucial for improving performance outcomes and promoting equal participation in mathematics education and careers.



\section{Origins and Longevity of Math Stereotypes}
The longevity of math stereotypes and their impact on women's participation in mathematics is a complex issue, deeply rooted in cultural beliefs and reinforced by various socialization agents.
\subsection{Gender-Role Congruence}
The belief that women are intrinsically less talented in mathematics or that they face more challenges in technological and mathematical fields is a significant aspect of the female gender role in Western culture (\cite{Tsui2007}). Women who internalize this cultural belief often have lower confidence in their arithmetic skills compared to their abilities in literature studies (\cite{Plante2009}).

Gender norms significantly influence personal beliefs and perceptions about the importance of various fields, ultimately discouraging women's interest in mathematics and the sciences. Contrary to the stereotypes of scientists and engineers, women often express a stronger interest in interpersonal relationships, helping others, and nurturing roles. In many descriptions, professionals are portrayed as lonely, overworked individuals who lack social and family time. This portrayal creates a disconnect for women, making math-related careers less appealing due to a mismatch between their evolving self-image and the perceived reality of STEM professions. The "image problem" (\cite{eccles1989}) associated with STEM fields contributes to this discrepancy, as these careers are often viewed as isolating and incompatible with a balanced lifestyle.

A lack of reliable information and effective reinforcement further exacerbates the problem. As previously stated, parents and teachers often provide inaccurate information about the actual mathematical skills needed for various careers and rarely present strong arguments for studying mathematics that is not required. The issue is further complicated by the lack of prominent female role models, which diminishes motivation for women to pursue advanced mathematics or related disciplines. Traditional gender roles can undermine women's confidence and appreciation for math, resulting in fewer female role models and discouraging them from entering STEM fields. This creates a vicious cycle that perpetuates gender stereotypes and reinforces negative perceptions of STEM careers across generations (\cite{eccles1989}).
\subsection{Parental Influences}
Parents' gender-stereotyped views often resist change, making them significant and frequently unintentional contributors to the formation of gender-differentiated values and self-perceptions (\cite{Gunderson2012}), even when confronted with their children's actual abilities. 

Parents of daughters often report that their child needs to put in more effort in math and perceive the subject as more challenging for girls. However, daughters frequently perform as well as or even better than sons in mathematics. A common bias contributes to this perception: parents are more likely to attribute boys' success in math to natural talent, while they tend to credit their daughters' achievements to hard work or diligence. Even when both genders receive equal grades, parents often view their sons as having greater mathematical ability because of this subtle yet powerful bias. Despite good intentions, this "effort versus talent" mindset suggests a lack of innate aptitude in girls, which undermines their self-efficacy and reinforces the idea that mathematical ability is inherently a male quality.

Parental confidence in their child's mathematical abilities significantly influences the child's values and self-perception, as shown by \textcite{Jacobs1991}. When parents perceive mathematics as particularly challenging for their daughters, it can lead to a decrease in the girls' confidence in the subject. This occurs because parents often frame academic performance within a context that may discourage female students. 

Moreover, parental guidance plays a crucial role in high school course selection. Parents tend to be less supportive of their daughters enrolling in advanced mathematics courses, often due to ingrained gender roles and a lack of exposure to non-traditional role models. As a result, well-meaning parents may unintentionally reinforce stereotypes, a phenomenon known as the "parental blind spot."

\subsection{Teacher Influences}
Teachers often reinforce rather than challenge gender-role stereotypes due to the practical demands of classroom management (\cite{Gunderson2012}). Through their interactions and the dynamics within the classroom, they unintentionally uphold these views, as they tend to react to situations rather than proactively address them. 

Subtle discrimination still exists in classrooms, although there has been a claimed decline in overt sexism. Male students are more likely to dominate interactions with their teachers and receive more public criticism, as was seen in reasearch by \textcite{VanLooRydell2014}. In many classrooms, a small number of male pupils often control class discussions, leaving the majority of students—including most females—as passive nonparticipants. This dynamic is largely due to teachers' unintentional responses to gender-role stereotyped behaviors; for example, males may be more assertive or exhibit greater familiarity with math and scientific equipment outside of the classroom. The resulting lack of essential interaction, feedback, and confidence-building opportunities for this "silent majority" of students—many of whom are girls—directly affects their motivation and perseverance in mathematics (\cite{Smith22}).

Research identifying "girl-friendly" or "superb" classrooms—those where girls exhibit more positive attitudes toward math—reveals several unique characteristics. These classrooms tend to be well-organized, have fewer compliments and criticisms, and operate more professionally. Teachers in these environments closely monitor interactions to ensure that every student participates. 

Such settings promote fewer public drills, less competition, and more one-on-one contact between teachers and students, all of which are beneficial for girls. In these "outstanding" classrooms, educators employ individualized or cooperative learning methods, avoid competitive motivational techniques, integrate experiential learning with real-world applications, and create dynamic and flexible learning environments. They also provide clear career counseling and ensure that all students engage in hands-on activities and computer use, preventing a few (primarily male) students from dominating these resources.
The results suggest that learning environments characterized by social comparison, competitiveness, lengthy public drills, and dominance by a small number of students are generally ineffective in motivating most female students to pursue science and math. In contrast, inclusive, collaborative, and hands-on environments are more beneficial for female students, minority students, and male students who struggle academically. It can be challenging for even well-intentioned instructors to ensure equitable participation due to structural issues related to classroom management and resource allocation, rather than solely individual teacher bias.

\section{Consequences and Broader Implications}
The lasting impact of math stereotypes and stereotype threat can significantly affect not only immediate academic performance but also long-term academic and career paths for women.
\subsection{Impact on Long-Term Academic and Career Choices}

This self-selection out of STEM fields is driven more by psychological discomfort than by a lack of ability or initial interest. The broader societal implications are significant, as women may feel pressured to conform or limit themselves, potentially sacrificing their full potential and the diverse perspectives that could enhance STEM fields. This academic disengagement effectively narrows their future career options, often before they even enter college (\cite{GadassiGati2009}). Furthermore, mathematical underperformance has been found to correlate with an increased propensity for engaging in stereotypically feminine activities (such as fashion, astrology, or culinary arts) rather than gender-neutral ones. This phenomenon suggests a "safe harbor" effect: when confronted with a threatening intellectual environment and diminished performance in a non-traditional domain, women may gravitate towards traditionally "safer" disciplines where their gender is positively stereotyped and their competence is affirmed. This implies a broader societal cost, as women are subtly pressured towards conformity or self-limitation, potentially sacrificing their full potential and diverse perspectives that could enrich STEM fields. 
\subsection{Theoretical Implications Regarding the Malleability of the Self-Concept and Identity Formation}
The observed modifications indicate that the self-concept is a flexible construct, capable of being adjusted to maintain a positive self-image. This adaptability is best exemplified by identity bifurcation, an ego-protective response that allows individuals, particularly women, to handle multiple identity needs simultaneously. This process highlights the unfairness of compelling people to abandon cherished interests or characteristics to "fit in," while also enabling perseverance in valued domains.

Furthermore, these findings align with established theories of identity formation. Carol Gilligan's work emphasizes the relational aspect of women's identity, where self-knowledge is developed through connections with others. In contrast, Erik Erikson's approach suggests that men's identity formation often relates to personal achievements. When faced with challenges in traditionally gendered fields, men and women draw confidence from different sources and employ various coping strategies, as these theoretical perspectives help to clarify. Women in STEM face hidden psychological costs due to systemic pressures to conform to gendered standards, rather than merely personal choice. 

\section{Recommendations for Future Research and Intervention}
Addressing the ongoing underrepresentation of women in mathematics and STEM fields requires a multifaceted approach that includes strategies to reduce stereotype threat, foster inclusive environments, and promote both competence and confidence through targeted interventions.
\subsection{Strategies for Lifting Stereotype Threat and Attenuating its Consequences}
An important first step is to actively challenge and eliminate negative perceptions. Additionally, it is essential to introduce diverse role models. Research shows that role models who can successfully balance multiple aspects of their identities—such as managing both work and family obligations—tend to have a greater impact. By breaking down limiting assumptions and demonstrating that success in STEM can align with a more integrated identity, these individuals play a significant role in reducing stereotype threats.
It is essential to develop both competence and confidence concurrently, as mere belief is insufficient for individuals to achieve tasks beyond their capabilities.
Teachers should actively engage in career counseling to highlight the importance of science and math for future employment opportunities, while also clarifying any misconceptions about the mathematical skills needed for various careers. Since men and women often derive their confidence from different sources, interventions should focus on building strong self-efficacy beliefs. For women, this involves emphasizing positive encouragement from influential individuals, providing relationship support, and offering vicarious experiences. For men, reinforcing mastery experiences remains a critical approach. 
\subsection{Future Research Directions}
Future research should further explore the complex interplay of factors that influence women in mathematics. This involves examining how cognitive and motivational aspects of stereotype threat act as mediating factors. It is also essential to identify the specific elements of stereotypes that cause discomfort and lead to performance deficits. Research should investigate whether certain achievement-related stereotypes about women's mathematical abilities are the primary source of stereotype anxiety. Longitudinal studies are necessary to differentiate the roles of individual coping abilities and situational factors in shaping women's academic and career paths. Additionally, including measures of gender orientation in self-efficacy assessments could help clarify whether differences in the development of confidence are rooted in beliefs about gender rather than gender itself. Finally, it is important to examine the self-efficacy development of ethnic and racial minority students, as there may be similarities with women's experiences in male-dominated fields.

\section{Aims of the research}
Despite ongoing initiatives aimed at promoting gender equity in education, notable disparities in mathematics attitudes, self-efficacy, and engagement between boys and girls persist —particularly during middle school, a critical period for the development of students' academic identities. Research indicates that pervasive gender stereotypes, such as the notion that mathematics is better suited for boys, can erode girls' confidence, diminish their interest in the subject, and restrict their aspirations in STEM fields (\cite{EcclesWigfield2020,Gunderson2012}). Additionally, the lack of visible female role models in mathematics reinforces these stereotypes, depriving students of counterexamples that could challenge the perception of mathematics as a male-dominated areas. Given that middle school is pivotal for forming gendered beliefs about mathematics, interventions that tackle these stereotypes and present alternative narratives are particularly essential.
This study aims to design, implement, and evaluate a short-term, school-based intervention that introduces middle school students to both historical and contemporary women mathematicians. The seminar is intended to challenge traditional gender stereotypes and foster more positive attitudes and self-beliefs regarding mathematics. The specific objectives of the study are as follows:
\begin{enumerate}
    \item \textbf{Challenge Gender Stereotypes in Mathematics:} The intervention will address common misconceptions about math being more suitable for boys, with the goal of promoting equitable perceptions of ability and potential among all students.
    \item \textbf{Increase Awareness of Women Mathematicians:} Students will learn about the achievements and career paths of female mathematicians, providing role models that counter stereotypical narratives and broaden their understanding of who can succeed in mathematics.
    \item \textbf{Examine Impacts on Attitudes and Self-Beliefs:} The study will investigate changes in students' self-efficacy, interest, enjoyment, and sense of belonging in mathematics, paying particular attention to differences between girls and boys.
    \item \textbf{Evaluate Effects on Gender Gaps in Outcomes:} The research will explore whether exposure to counter-stereotypical role models reduces gender differences in affective outcomes related to mathematics, including confidence, attitudes, and engagement.
\end{enumerate}

Through these objectives, the research intends to provide evidence on the effectiveness of brief, stereotype-challenging interventions in middle school classrooms. The intervention is expected to yield several positive outcomes: girls are anticipated to demonstrate increased self-efficacy, enjoyment, and sense of belonging in mathematics; both girls and boys are expected to show reduced endorsement of traditional gender stereotypes; and although the primary focus is on attitudinal and motivational outcomes, enhanced engagement may also indirectly influence performance. Overall, the study seeks to contribute to a deeper understanding of how targeted educational interventions can promote gender equity in mathematics education during a critical stage of students' schooling.